\section{Succession Rules for the Early Replacement of the Government or Parts of It}\label{sec:succession_rules}

For many reasons, it can happen that the Bundestag and the Federal Government, or only parts of them,
have to be replaced early.
So that the country is never left without a government, there must be replacement rules for all conceivable cases.
The rules currently in force cannot be adopted without modification.
The proposed procedural changes are too fundamental for that.
For example, the Chancellor is now to be elected indirectly by the people and no longer by the Bundestag.
Therefore, the Bundestag cannot simply elect a new Chancellor.
That would only be an option as a temporary interim solution.

In the following, the conceivable cases are therefore gone through one by one:

\subsection{Replacement of the Chancellor Without a Change in the Legislature}\label{subsec:chancellor_replacement_without_legislative_change}
It can happen that the Chancellor has to be replaced early although nothing else has changed in the balance of power
in parliament.

The Chancellor could, e.g., become seriously ill or even die, resign of their own accord
or no longer be allowed to hold office because they have personally committed something so serious
that they may no longer be Chancellor.
In this case, no new parliamentary elections need to be held.
It is sufficient to elect the Chancellor anew.

The party from which the outgoing Chancellor comes must again present three chancellor candidates to the people for election.
Two of the three candidates may, but need not, be the same as in the original federal election.
In a single-party government, the candidate with the most votes automatically becomes Chancellor.
The Federal President must appoint them.

In a coalition government, the procedure from the \mvref{subsec:election-of-lead-candidates-by-the-people} applies.
The Chancellor can be determined by method a) or b) and requires confirmation by the coalition partners
in a secret ballot.

If this confirmation is not given, a new coalition can still be formed analogously to \cref{subsec:coalition_breakup_with_new_coalition}.
The only difference from \cref{subsec:coalition_breakup_with_new_coalition} is that option c) is naturally no longer
available.
Should the old chancellor's party be part of the new coalition, it enters this coalition with the result of the chancellor-candidate by-election.\footnote{The term \enquote{new election of the Chancellor} would not be correct, since whether any of the new chancellor candidates becomes Chancellor at all also depends on the election results of the other coalition partners.}

If forming a new coalition is not wanted or does not succeed, new Bundestag elections must be held.
Until the Federal President appoints a new Chancellor, the Vice Chancellor conducts the business of government on an acting basis.

\runinhead{Note:}
Even if it is clear from the outset that the existing governing coalition will break up over the missing Chancellor,
a chancellor-candidate by-election must still be held if the old chancellor's party takes part in a new coalition.
This applies even if it joins a new coalition not immediately after the old one has broken up,
but, for example, only after one or more further new coalitions have broken up.
The reason is that the party no longer has a chancellor candidate chosen by the people from three candidates.

The Federal President moderates this process.
The President orders a chancellor-candidate by-election or refrains from doing so if a coalition is formed without the old chancellor's party,
or orders new parliamentary elections as soon as it is clear that no successor coalition will be found.
In the case of a single-party government, which precisely does not depend on a coalition to govern, the Federal President
always orders a chancellor-candidate by-election if only the Chancellor drops out.

A single-party government, too, may have to be replaced as a whole, e.g.\ because the party falls out with itself.
This case is to be treated like the breakup of a coalition and to be distinguished from the replacement of the person of the Chancellor described here.

\subsection{Breakup of a Coalition Followed by the Formation of a New Coalition} \label{subsec:coalition_breakup_with_new_coalition}
It often happens that coalitions break up, and sometimes new ones come together afterwards.
Once a new coalition has come together, it must inform the Federal President so that the President orders chancellor-candidate by-elections
in the event that a party cannot present three chancellor candidates elected by the people.
See \cref{subsec:chancellor_replacement_without_legislative_change}.

Once it is ensured that all coalition parties enter the coalition with chancellor candidates elected by the people, the test of a
successful coalition formation is the election of the Chancellor by the coalition in accordance with the \mvref{subsec:election-of-lead-candidates-by-the-people}.
Unlike coalition formations directly after a federal election, however, there should additionally be the option for
the old Chancellor to remain in office if the coalition partners agree on this and the old Chancellor is available
for election.

This applies even if the old Chancellor does not belong to any party of the new coalition and would thus govern without a party of their own
behind them.
The purpose of this rule is to enable maximum stability of the executive.
If one is lucky enough to have a Chancellor who is experienced in office and respected as a person by all
parties, then they should be allowed to continue in office despite party disputes if they themselves
want this and the new coalition partners want it too.

The coalition partners thus have the choice between candidates determined by three different methods:
\begin{enumerate}[label=\alph*)]
  \item The candidate of the strongest coalition party who received the most votes.
  \item The candidate from all coalition parties who received the most votes.
  \item The old Chancellor.
\end{enumerate}

These three methods can yield three different candidates, or only two or one, because, e.g.,
the old chancellor's party is also the strongest party in the new coalition and the old Chancellor is also
the candidate who won the most votes overall in the last federal election.

Should no Chancellor be confirmed in the secret chancellor election (possibly in several rounds), the coalition
has not come about, and the Federal President orders new parliamentary elections as soon as it is clear that no other coalition
will come about and that the incumbent Chancellor will not form a minority government either.

\subsection{Breakup of a Coalition Without the Formation of a New Coalition}
If a coalition supporting the Chancellor breaks up and no new one is formed, this does not necessarily mean that
the Chancellor would have to resign because of it.
The Chancellor can continue to exercise the executive office even without a majority of their own in the legislature, as long as they are not explicitly removed.
If they need the approval of parliament for their plans, they simply have to find suitable majorities for each plan.

If a governing coalition breaks up, no new one is found to replace it and the old Chancellor does not want to lead a minority government
either, the Federal President must order new parliamentary elections.
The old Chancellor must inform the Federal President of this.

\subsection{Replacement of Individual Ministers}
The procedure for replacing federal ministers remains as it is: the Federal President appoints or dismisses them on the proposal of the Chancellor.

\subsection{Replacement of Individual Members of Parliament} \label{subsec:replacement_of_individual_mps}
If individual members of parliament have to be replaced, list candidates are simply succeeded by the next candidates on the list.
For directly elected members, a by-election must be held in the assigned constituency.
